%\newpage

{\centering
\section{Conclusioni}
}
\medskip
{\color{gray}\hrule}
\medskip
In questa esperienza si sono utilizzate le transizioni atomiche di una lampada ai vapori di Cadmio 
per stimare il magnetone di Bohr osservando l'effetto Zeeman.

Per effettuare le misure di campo magnetico, si è calibrato un ferromagnete, mettendo in relazione il campo da esso generato con la corrente che scorre attraverso esso: 
$$B(I) = (63.9\pm1.2)\text{mT/A}\cdot I + (12\pm5)\text{mT}$$
questa relazione è lineare, e si è inoltre verificato che la risoluzione dell'apparato non permette di distinguere la discesa e la salita del ciclo di isteresi del ferromagnete.

Si è poi utilizzato un banco ottico per distinguere con tecniche interferometriche le diverse transizioni atomiche prodotte dall'effetto Zeeman e stimarne la differenza in energia: questa, per la teoria, è direttamente proporzionale al campo magnetico in cui la lampada è immersa e da questa costante di proporzionalità si sono ottenute diverse stime del magnetone di Bohr $\mu_B$.\\
In conclusione, le tre stime estratte e la loro media pesata
$$\langle\mu_B\rangle = (9.6\pm0.2)\cdot10^{-24}\text{JT}^{-1}$$
sono risultate compatibili con il valore teorico di riferimento.